\documentclass[../main.tex]{subfiles}

\begin{document}

\section{Conclusions}
\justify

Taken as a whole, the experimental results are encouraging. The simulated drone reached an episode-level success rate of 89.11\,\% (and a per-obstacle success rate of 95.60\,\%) across the two density conditions, and that performance carried over to the real forest with an episode-level transfer efficiency of 81.11\,\% and a per-obstacle transfer efficiency of 84.02\,\%. Neither figure is 100\,\%, which means the system is clearly not production-ready yet, but the gap between the simulated and real-world numbers is small enough -- and the absolute numbers are high enough across 200 evaluation episodes -- that the results are statistically meaningful rather than noise.

The numbers become more impressive once the platform constraints are taken into account. The drone was equipped with only a single \$20 monocular camera, no LiDAR or stereo, and was carrying that camera on a full-size Holybro S500 V2 frame whose footprint is much larger than the small agile racers usually used in vision-based avoidance studies. The camera itself runs at a narrow $62^{\circ}$ field of view, which is unusually restrictive for forest flight and clips obstacles from the periphery sooner than the wide-angle lenses used in comparable work. That this configuration still reaches above 80\,\% transfer is the central empirical result of this thesis.

The fact that those numbers were obtained with such a minimal sensor setup also points forward. The camera is small enough that the airframe could likely be shrunk significantly without losing perception capacity, and either widening the field of view or adding a second camera to recover lost peripheral information would be a natural next step that could plausibly push the transfer numbers higher than the ones reported here.

\end{document}
